
%% bare_jrnl_compsoc.tex
%% V1.4a
%% 2014/09/17
%% by Michael Shell
%% See:
%% http://www.michaelshell.org/
%% for current contact information.
%%
%% This is a skeleton file demonstrating the use of IEEEtran.cls
%% (requires IEEEtran.cls version 1.8a or later) with an IEEE
%% Computer Society journal paper.
%%
%% Support sites:
%% http://www.michaelshell.org/tex/ieeetran/
%% http://www.ctan.org/tex-archive/macros/latex/contrib/IEEEtran/
%% and
%% http://www.ieee.org/

%%*************************************************************************
%% Legal Notice:
%% This code is offered as-is without any warranty either expressed or
%% implied; without even the implied warranty of MERCHANTABILITY or
%% FITNESS FOR A PARTICULAR PURPOSE! 
%% User assumes all risk.
%% In no event shall IEEE or any contributor to this code be liable for
%% any damages or losses, including, but not limited to, incidental,
%% consequential, or any other damages, resulting from the use or misuse
%% of any information contained here.
%%
%% All comments are the opinions of their respective authors and are not
%% necessarily endorsed by the IEEE.
%%
%% This work is distributed under the LaTeX Project Public License (LPPL)
%% ( http://www.latex-project.org/ ) version 1.3, and may be freely used,
%% distributed and modified. A copy of the LPPL, version 1.3, is included
%% in the base LaTeX documentation of all distributions of LaTeX released
%% 2003/12/01 or later.
%% Retain all contribution notices and credits.
%% ** Modified files should be clearly indicated as such, including  **
%% ** renaming them and changing author support contact information. **
%%
%% File list of work: IEEEtran.cls, IEEEtran_HOWTO.pdf, bare_adv.tex,
%%                    bare_conf.tex, bare_jrnl.tex, bare_conf_compsoc.tex,
%%                    bare_jrnl_compsoc.tex, bare_jrnl_transmag.tex
%%*************************************************************************


% *** Authors should verify (and, if needed, correct) their LaTeX system  ***
% *** with the testflow diagnostic prior to trusting their LaTeX platform ***
% *** with production work. IEEE's font choices and paper sizes can       ***
% *** trigger bugs that do not appear when using other class files.       ***                          ***
% The testflow support page is at:
% http://www.michaelshell.org/tex/testflow/


\documentclass[10pt,journal]{IEEEtran}
%
% If IEEEtran.cls has not been installed into the LaTeX system files,
% manually specify the path to it like:
% \documentclass[10pt,journal,compsoc]{../sty/IEEEtran}


\bibliographystyle{IEEEtran}
%\bibliographystyle{natdin}

% Some very useful LaTeX packages include:
% (uncomment the ones you want to load)


% *** MISC UTILITY PACKAGES ***
%
%\usepackage{ifpdf}
% Heiko Oberdiek's ifpdf.sty is very useful if you need conditional
% compilation based on whether the output is pdf or dvi.
% usage:
% \ifpdf
%   % pdf code
% \else
%   % dvi code
% \fi
% The latest version of ifpdf.sty can be obtained from:
% http://www.ctan.org/tex-archive/macros/latex/contrib/oberdiek/
% Also, note that IEEEtran.cls V1.7 and later provides a builtin
% \ifCLASSINFOpdf conditional that works the same way.
% When switching from latex to pdflatex and vice-versa, the compiler may
% have to be run twice to clear warning/error messages.


\newcommand{\ezio}[1]{\noindent\textcolor{red}{[EZIO] #1}}



% *** CITATION PACKAGES ***
%
\ifCLASSOPTIONcompsoc
  % IEEE Computer Society needs nocompress option
  % requires cite.sty v4.0 or later (November 2003)
  \usepackage[nocompress]{cite}
\else
  % normal IEEE
  \usepackage{cite}
\fi
% cite.sty was written by Donald Arseneau
% V1.6 and later of IEEEtran pre-defines the format of the cite.sty package
% \cite{} output to follow that of IEEE. Loading the cite package will
% result in citation numbers being automatically sorted and properly
% "compressed/ranged". e.g., [1], [9], [2], [7], [5], [6] without using
% cite.sty will become [1], [2], [5]--[7], [9] using cite.sty. cite.sty's
% \cite will automatically add leading space, if needed. Use cite.sty's
% noadjust option (cite.sty V3.8 and later) if you want to turn this off
% such as if a citation ever needs to be enclosed in parenthesis.
% cite.sty is already installed on most LaTeX systems. Be sure and use
% version 5.0 (2009-03-20) and later if using hyperref.sty.
% The latest version can be obtained at:
% http://www.ctan.org/tex-archive/macros/latex/contrib/cite/
% The documentation is contained in the cite.sty file itself.
%
% Note that some packages require special options to format as the Computer
% Society requires. In particular, Computer Society  papers do not use
% compressed citation ranges as is done in typical IEEE papers
% (e.g., [1]-[4]). Instead, they list every citation separately in order
% (e.g., [1], [2], [3], [4]). To get the latter we need to load the cite
% package with the nocompress option which is supported by cite.sty v4.0
% and later. Note also the use of a CLASSOPTION conditional provided by
% IEEEtran.cls V1.7 and later.





% *** GRAPHICS RELATED PACKAGES ***
%
\ifCLASSINFOpdf
  % \usepackage[pdftex]{graphicx}
  % declare the path(s) where your graphic files are
  % \graphicspath{{../pdf/}{../jpeg/}}
  % and their extensions so you won't have to specify these with
  % every instance of \includegraphics
  % \DeclareGraphicsExtensions{.pdf,.jpeg,.png}
\else
  % or other class option (dvipsone, dvipdf, if not using dvips). graphicx
  % will default to the driver specified in the system graphics.cfg if no
  % driver is specified.
  % \usepackage[dvips]{graphicx}
  % declare the path(s) where your graphic files are
  % \graphicspath{{../eps/}}
  % and their extensions so you won't have to specify these with
  % every instance of \includegraphics
  % \DeclareGraphicsExtensions{.eps}
\fi
% graphicx was written by David Carlisle and Sebastian Rahtz. It is
% required if you want graphics, photos, etc. graphicx.sty is already
% installed on most LaTeX systems. The latest version and documentation
% can be obtained at: 
% http://www.ctan.org/tex-archive/macros/latex/required/graphics/
% Another good source of documentation is "Using Imported Graphics in
% LaTeX2e" by Keith Reckdahl which can be found at:
% http://www.ctan.org/tex-archive/info/epslatex/
%
% latex, and pdflatex in dvi mode, support graphics in encapsulated
% postscript (.eps) format. pdflatex in pdf mode supports graphics
% in .pdf, .jpeg, .png and .mps (metapost) formats. Users should ensure
% that all non-photo figures use a vector format (.eps, .pdf, .mps) and
% not a bitmapped formats (.jpeg, .png). IEEE frowns on bitmapped formats
% which can result in "jaggedy"/blurry rendering of lines and letters as
% well as large increases in file sizes.
%
% You can find documentation about the pdfTeX application at:
% http://www.tug.org/applications/pdftex






% *** MATH PACKAGES ***
%
%\usepackage[cmex10]{amsmath}
% A popular package from the American Mathematical Society that provides
% many useful and powerful commands for dealing with mathematics. If using
% it, be sure to load this package with the cmex10 option to ensure that
% only type 1 fonts will utilized at all point sizes. Without this option,
% it is possible that some math symbols, particularly those within
% footnotes, will be rendered in bitmap form which will result in a
% document that can not be IEEE Xplore compliant!
%
% Also, note that the amsmath package sets \interdisplaylinepenalty to 10000
% thus preventing page breaks from occurring within multiline equations. Use:
%\interdisplaylinepenalty=2500
% after loading amsmath to restore such page breaks as IEEEtran.cls normally
% does. amsmath.sty is already installed on most LaTeX systems. The latest
% version and documentation can be obtained at:
% http://www.ctan.org/tex-archive/macros/latex/required/amslatex/math/





% *** SPECIALIZED LIST PACKAGES ***
%
%\usepackage{algorithmic}
% algorithmic.sty was written by Peter Williams and Rogerio Brito.
% This package provides an algorithmic environment fo describing algorithms.
% You can use the algorithmic environment in-text or within a figure
% environment to provide for a floating algorithm. Do NOT use the algorithm
% floating environment provided by algorithm.sty (by the same authors) or
% algorithm2e.sty (by Christophe Fiorio) as IEEE does not use dedicated
% algorithm float types and packages that provide these will not provide
% correct IEEE style captions. The latest version and documentation of
% algorithmic.sty can be obtained at:
% http://www.ctan.org/tex-archive/macros/latex/contrib/algorithms/
% There is also a support site at:
% http://algorithms.berlios.de/index.html
% Also of interest may be the (relatively newer and more customizable)
% algorithmicx.sty package by Szasz Janos:
% http://www.ctan.org/tex-archive/macros/latex/contrib/algorithmicx/




% *** ALIGNMENT PACKAGES ***
%
%\usepackage{array}
% Frank Mittelbach's and David Carlisle's array.sty patches and improves
% the standard LaTeX2e array and tabular environments to provide better
% appearance and additional user controls. As the default LaTeX2e table
% generation code is lacking to the point of almost being broken with
% respect to the quality of the end results, all users are strongly
% advised to use an enhanced (at the very least that provided by array.sty)
% set of table tools. array.sty is already installed on most systems. The
% latest version and documentation can be obtained at:
% http://www.ctan.org/tex-archive/macros/latex/required/tools/


% IEEEtran contains the IEEEeqnarray family of commands that can be used to
% generate multiline equations as well as matrices, tables, etc., of high
% quality.




% *** SUBFIGURE PACKAGES ***
%\ifCLASSOPTIONcompsoc
%  \usepackage[caption=false,font=footnotesize,labelfont=sf,textfont=sf]{subfig}
%\else
%  \usepackage[caption=false,font=footnotesize]{subfig}
%\fi
% subfig.sty, written by Steven Douglas Cochran, is the modern replacement
% for subfigure.sty, the latter of which is no longer maintained and is
% incompatible with some LaTeX packages including fixltx2e. However,
% subfig.sty requires and automatically loads Axel Sommerfeldt's caption.sty
% which will override IEEEtran.cls' handling of captions and this will result
% in non-IEEE style figure/table captions. To prevent this problem, be sure
% and invoke subfig.sty's "caption=false" package option (available since
% subfig.sty version 1.3, 2005/06/28) as this is will preserve IEEEtran.cls
% handling of captions.
% Note that the Computer Society format requires a sans serif font rather
% than the serif font used in traditional IEEE formatting and thus the need
% to invoke different subfig.sty package options depending on whether
% compsoc mode has been enabled.
%
% The latest version and documentation of subfig.sty can be obtained at:
% http://www.ctan.org/tex-archive/macros/latex/contrib/subfig/




% *** FLOAT PACKAGES ***
%
%\usepackage{fixltx2e}
% fixltx2e, the successor to the earlier fix2col.sty, was written by
% Frank Mittelbach and David Carlisle. This package corrects a few problems
% in the LaTeX2e kernel, the most notable of which is that in current
% LaTeX2e releases, the ordering of single and double column floats is not
% guaranteed to be preserved. Thus, an unpatched LaTeX2e can allow a
% single column figure to be placed prior to an earlier double column
% figure. The latest version and documentation can be found at:
% http://www.ctan.org/tex-archive/macros/latex/base/


%\usepackage{stfloats}
% stfloats.sty was written by Sigitas Tolusis. This package gives LaTeX2e
% the ability to do double column floats at the bottom of the page as well
% as the top. (e.g., "\begin{figure*}[!b]" is not normally possible in
% LaTeX2e). It also provides a command:
%\fnbelowfloat
% to enable the placement of footnotes below bottom floats (the standard
% LaTeX2e kernel puts them above bottom floats). This is an invasive package
% which rewrites many portions of the LaTeX2e float routines. It may not work
% with other packages that modify the LaTeX2e float routines. The latest
% version and documentation can be obtained at:
% http://www.ctan.org/tex-archive/macros/latex/contrib/sttools/
% Do not use the stfloats baselinefloat ability as IEEE does not allow
% \baselineskip to stretch. Authors submitting work to the IEEE should note
% that IEEE rarely uses double column equations and that authors should try
% to avoid such use. Do not be tempted to use the cuted.sty or midfloat.sty
% packages (also by Sigitas Tolusis) as IEEE does not format its papers in
% such ways.
% Do not attempt to use stfloats with fixltx2e as they are incompatible.
% Instead, use Morten Hogholm'a dblfloatfix which combines the features
% of both fixltx2e and stfloats:
%
% \usepackage{dblfloatfix}
% The latest version can be found at:
% http://www.ctan.org/tex-archive/macros/latex/contrib/dblfloatfix/




%\ifCLASSOPTIONcaptionsoff
%  \usepackage[nomarkers]{endfloat}
% \let\MYoriglatexcaption\caption
% \renewcommand{\caption}[2][\relax]{\MYoriglatexcaption[#2]{#2}}
%\fi
% endfloat.sty was written by James Darrell McCauley, Jeff Goldberg and 
% Axel Sommerfeldt. This package may be useful when used in conjunction with 
% IEEEtran.cls'  captionsoff option. Some IEEE journals/societies require that
% submissions have lists of figures/tables at the end of the paper and that
% figures/tables without any captions are placed on a page by themselves at
% the end of the document. If needed, the draftcls IEEEtran class option or
% \CLASSINPUTbaselinestretch interface can be used to increase the line
% spacing as well. Be sure and use the nomarkers option of endfloat to
% prevent endfloat from "marking" where the figures would have been placed
% in the text. The two hack lines of code above are a slight modification of
% that suggested by in the endfloat docs (section 8.4.1) to ensure that
% the full captions always appear in the list of figures/tables - even if
% the user used the short optional argument of \caption[]{}.
% IEEE papers do not typically make use of \caption[]'s optional argument,
% so this should not be an issue. A similar trick can be used to disable
% captions of packages such as subfig.sty that lack options to turn off
% the subcaptions:
% For subfig.sty:
% \let\MYorigsubfloat\subfloat
% \renewcommand{\subfloat}[2][\relax]{\MYorigsubfloat[]{#2}}
% However, the above trick will not work if both optional arguments of
% the \subfloat command are used. Furthermore, there needs to be a
% description of each subfigure *somewhere* and endfloat does not add
% subfigure captions to its list of figures. Thus, the best approach is to
% avoid the use of subfigure captions (many IEEE journals avoid them anyway)
% and instead reference/explain all the subfigures within the main caption.
% The latest version of endfloat.sty and its documentation can obtained at:
% http://www.ctan.org/tex-archive/macros/latex/contrib/endfloat/
%
% The IEEEtran \ifCLASSOPTIONcaptionsoff conditional can also be used
% later in the document, say, to conditionally put the References on a 
% page by themselves.




% *** PDF, URL AND HYPERLINK PACKAGES ***
%
%\usepackage{url}
% url.sty was written by Donald Arseneau. It provides better support for
% handling and breaking URLs. url.sty is already installed on most LaTeX
% systems. The latest version and documentation can be obtained at:
% http://www.ctan.org/tex-archive/macros/latex/contrib/url/
% Basically, \url{my_url_here}.





% *** Do not adjust lengths that control margins, column widths, etc. ***
% *** Do not use packages that alter fonts (such as pslatex).         ***
% There should be no need to do such things with IEEEtran.cls V1.6 and later.
% (Unless specifically asked to do so by the journal or conference you plan
% to submit to, of course. )

\usepackage{times}
\usepackage{epsfig}
\usepackage{ifpdf}
\usepackage{stackrel}

\usepackage{array}
\usepackage{amsmath,amssymb,amscd,amstext}
% *** MATH PACKAGES ***
\usepackage{mathrsfs}
\usepackage{dsfont}
\usepackage{amsbsy}
\usepackage{stmaryrd}

\usepackage{subfigure}

%\usepackage{tocloft}
\usepackage{setspace}
%\usepackage{proof}
\usepackage{fancyhdr} 
\usepackage{hyperref}
\usepackage{lastpage}
\usepackage{pslatex}
\usepackage{amsthm}

\usepackage[verbose]{placeins}

\usepackage{amsmath}

%\usepackage{biblatex}
%\usepackage{csquotes}
%\bibliography{nfm} 

\usepackage[displaymath,textmath,sections,graphics,floats,graphicx,color]{preview}
%\newcommand{\tuple}[1]{\langle #1\rangle}
\def  \Nat      {\mathbb{N}}   % naturals
\def  \Int      {\mathbb{Z}}     % integers
\def  \Real     {\mathbb{R}}  % reals
\def  \Bool     {\mathbb{B}}
\def  \Complex  {\mathbb{C}}
\def  \T        {\mathbb{T}}
\def  \Domain   {\mathbb{D}}

\def  \om       {\omega}
\def  \hf       {\hat{f}}
\def  \hx       {\hat{x}}
\def  \ft#1     {\footnote{#1} }
\def  \f        {\varphi}
\def  \ni       {\noindent}
\def  \t        {\tau}
\def  \s        {\sigma}

\renewcommand{\i}{\mathbf{i}}
%\newcommand{\ev}{\Diamond}
%\renewcommand{\next}{\bigcirc}
%\newcommand{\always}{\square}
%\newcommand{\untz}{~\mathcal{U}~}   % no argument
%\newcommand{\unti}[1]{~\mathcal{U}_{#1}}
%\newcommand{\untii}[2]{~\mathcal{U}_{[#1,#2]}}
\newcommand{\sncz}{~\mathcal{S}~}   % no argument
\newcommand{\snci}[1]{~\mathcal{S}_{#1}}
\newcommand{\sncii}[2]{~\mathcal{S}_{[#1,#2]}}
\newcommand{\hz}{\mbox{Hz}}
\newcommand{\pitch}{\mathop{\text{pitch}}}


\newcommand{\lts}[1]{\ensuremath{\stackrel{#1}{\rightarrow}}}


\newcommand{\until}[1]{\ensuremath{\mathbf{U}_{#1}}}
\newcommand{\nxt}[1]{\ensuremath{\mathbf{X}_{#1}}}
\newcommand{\always}[1]{\ensuremath{\mathbf{G}_{#1}}}
\newcommand{\eventually}[1]{\ensuremath{\mathbf{F}_{#1}}}
\newcommand{\true}{\mathtt{true}}
\newcommand{\false}{\mathtt{false}}

\usepackage{tikz}
\usepackage{xcolor}
\usepackage{pdflscape}

\usepackage{algorithm, algpseudocode}



%\newcommand{\infrule}[1]{\scriptstyle\it{#1}}


\usetikzlibrary{arrows,positioning,automata}


% % % Thm environments:
\newtheorem{remark}{\bf Remark}[section]
\newtheorem{proposition}{\bf Proposition}[section]
\newtheorem{ass}{\bf Assumption}[section]
\newtheorem{definition}{\bf Definition}[section]
\newtheorem{example}{\bf Example}[section]
\newtheorem{theorem}{\bf Theorem}[section]
\newtheorem{problem}{\bf Problem}[section]


\renewcommand{\theequation}{\thesection.\arabic{equation}}
\newcommand{\red}{\textcolor{red}}


\renewcommand{\vec}[1]{\mathbf{#1}}
\renewcommand{\cal}[1]{\mathcal{#1}}
\newcommand{\bb}[1]{\mathbb{#1}}
\renewcommand{\phi}{\varphi}


%\newcommand{\change}[1]{\textcolor{red}{#1}}
%replace with the following one to remove the red color
\newcommand{\change}[1]{#1}

% \pr(A \cond B) denotes a conditional probability
\newcommand{\cond}{\mathrel{|}}
% \pr(A \joint B) denotes a joint probability
\newcommand{\joint}{\cap}



% correct bad hyphenation here
\hyphenation{op-tical net-works semi-conduc-tor}


\begin{document}
% GENERIC MATH MACROS
%\newcommand{\eventually}{\Diamond}
%\newcommand{\henceforth}{\Box}
\newcommand{\tuple}[1]{\langle #1\rangle}
\newcommand{\sat}{\models}
% \pr(X) = probability of X.  it's probably a good idea to use same font
% for \pr and \Peek.
\newcommand{\pr}{{\it P}}
\newcommand{\Peek}{{\it Peek}}
\newcommand{\comment}[1]{}

% \fullOnly{material that appears in tech report but not conference paper}
\newcommand{\fullOnly}[1]{}
%\newcommand{\fullOnly}[1]{#1}

% \confOnly{material that appears in conference paper but not tech report}
%\newcommand{\confOnly}[1]{}
\newcommand{\confOnly}[1]{#1}

% DOCUMENT-SPECIFIC MACROS
% observation symbol representing a gap
\newcommand{\gap}{{\it gap}}
% random variable representing the result of a peek operation.
% use regular font not \mathbf, i.e., assume a single variable,
% otherwise it looks odd, because X, O, and S are single variables.
\newcommand{\pkvar}{Q}
% meta-variable representing a value returned by a peek operation,
% i.e., a value of \pkvar.
\newcommand{\pk}{q}

% \sampled{O} = sampled version of trace O
%\newcommand{\samp}[1]{\check{#1}}
%\newcommand{\samp}[1]{{\rm samp}(#1)}

%%% Local Variables: 
%%% mode: latex
%%% TeX-master: "main"
%%% End: 

%
% paper title
% Titles are generally capitalized except for words such as a, an, and, as,
% at, but, by, for, in, nor, of, on, or, the, to and up, which are usually
% not capitalized unless they are the first or last word of the title.
% Linebreaks \\ can be used within to get better formatting as desired.
% Do not put math or special symbols in the title.
\title{Matrix inversion using CUDA}
%
%
% author names and IEEE memberships
% note positions of commas and nonbreaking spaces ( ~ ) LaTeX will not break
% a structure at a ~ so this keeps an author's name from being broken across
% two lines.
% use \thanks{} to gain access to the first footnote area
% a separate \thanks must be used for each paragraph as LaTeX2e's \thanks
% was not built to handle multiple paragraphs
%
%
%\IEEEcompsocitemizethanks is a special \thanks that produces the bulleted
% lists the Computer Society journals use for "first footnote" author
% affiliations. Use \IEEEcompsocthanksitem which works much like \item
% for each affiliation group. When not in compsoc mode,
% \IEEEcompsocitemizethanks becomes like \thanks and
% \IEEEcompsocthanksitem becomes a line break with idention. This
% facilitates dual compilation, although admittedly the differences in the
% desired content of \author between the different types of papers makes a
% one-size-fits-all approach a daunting prospect. For instance, compsoc 
% journal papers have the author affiliations above the "Manuscript
% received ..."  text while in non-compsoc journals this is reversed. Sigh.


\author{Mathias Lechner (1225134),~\IEEEmembership{}
     Jakob L\"ovhall (1527836),~\IEEEmembership{}
        Patrick Kn\"obel~\IEEEmembership{}% <-this % stops a space
% <-this % stops an unwanted space
%\thanks{Manuscript received April 19, 2005; revised September 17, 2014.}
}





%\author{Michael~Shell,~\IEEEmembership{Member,~IEEE,}
%   %  John~Doe,~\IEEEmembership{Fellow,~OSA,}
%        and~Jane~Doe,~\IEEEmembership{Life~Fellow,~IEEE}% <-this % stops a space
%\IEEEcompsocitemizethanks{\IEEEcompsocthanksitem M. Shell is with the Department
%of Electrical and Computer Engineering, Georgia Institute of Technology, Atlanta,
%GA, 30332.\protect\\
% %note need leading \protect in front of \\ to get a newline within \thanks as
% %\\ is fragile and will error, could use \hfil\break instead.
%E-mail: see http://www.michaelshell.org/contact.html
%\IEEEcompsocthanksitem J. Doe and J. Doe are with Anonymous University.}% <-this % stops an unwanted space
%\thanks{Manuscript received April 19, 2005; revised September 17, 2014.}}

% note the % following the last \IEEEmembership and also \thanks - 
% these prevent an unwanted space from occurring between the last author name
% and the end of the author line. i.e., if you had this:
% 
% \author{....lastname \thanks{...} \thanks{...} }
%                     ^------------^------------^----Do not want these spaces!
%
% a space would be appended to the last name and could cause every name on that
% line to be shifted left slightly. This is one of those "LaTeX things". For
% instance, "\textbf{A} \textbf{B}" will typeset as "A B" not "AB". To get
% "AB" then you have to do: "\textbf{A}\textbf{B}"
% \thanks is no different in this regard, so shield the last } of each \thanks
% that ends a line with a % and do not let a space in before the next \thanks.
% Spaces after \IEEEmembership other than the last one are OK (and needed) as
% you are supposed to have spaces between the names. For what it is worth,
% this is a minor point as most people would not even notice if the said evil
% space somehow managed to creep in.



% The paper headers
%\markboth{Journal of \LaTeX\ Class Files,~Vol.~13, No.~9, September~2014}%
%{Shell \MakeLowercase{\textit{et al.}}: Bare Demo of IEEEtran.cls for Computer Society Journals}
% The only time the second header will appear is for the odd numbered pages
% after the title page when using the twoside option.
% 
% *** Note that you probably will NOT want to include the author's ***
% *** name in the headers of peer review papers.                   ***
% You can use \ifCLASSOPTIONpeerreview for conditional compilation here if
% you desire.



% The publisher's ID mark at the bottom of the page is less important with
% Computer Society journal papers as those publications place the marks
% outside of the main text columns and, therefore, unlike regular IEEE
% journals, the available text space is not reduced by their presence.
% If you want to put a publisher's ID mark on the page you can do it like
% this:
%\IEEEpubid{0000--0000/00\$00.00~\copyright~2014 IEEE}
% or like this to get the Computer Society new two part style.
%\IEEEpubid{\makebox[\columnwidth]{\hfill 0000--0000/00/\$00.00~\copyright~2014 IEEE}%
%\hspace{\columnsep}\makebox[\columnwidth]{Published by the IEEE Computer Society\hfill}}
% Remember, if you use this you must call \IEEEpubidadjcol in the second
% column for its text to clear the IEEEpubid mark (Computer Society jorunal
% papers don't need this extra clearance.)



% use for special paper notices
%\IEEEspecialpapernotice{(Invited Paper)}



% for Computer Society papers, we must declare the abstract and index terms
% PRIOR to the title within the \IEEEtitleabstractindextext IEEEtran
% command as these need to go into the title area created by \maketitle.
% As a general rule, do not put math, special symbols or citations
% in the abstract or keywords.
\IEEEtitleabstractindextext{%
% !TEX root =  root.tex
\begin{abstract}

\end{abstract}

}

% Note that keywords are not normally used for peerreview papers.
%\begin{IEEEkeywords}
%Computer Society, IEEEtran, journal, \LaTeX, paper, template.
%\end{IEEEkeywords}}


% make the title area
\maketitle


% To allow for easy dual compilation without having to reenter the
% abstract/keywords data, the \IEEEtitleabstractindextext text will
% not be used in maketitle, but will appear (i.e., to be "transported")
% here as \IEEEdisplaynontitleabstractindextext when the compsoc 
% or transmag modes are not selected <OR> if conference mode is selected 
% - because all conference papers position the abstract like regular
% papers do.
\IEEEdisplaynontitleabstractindextext
% \IEEEdisplaynontitleabstractindextext has no effect when using
% compsoc or transmag under a non-conference mode.



% For peer review papers, you can put extra information on the cover
% page as needed:
% \ifCLASSOPTIONpeerreview
% \begin{center} \bfseries EDICS Category: 3-BBND \end{center}
% \fi
%
% For peerreview papers, this IEEEtran command inserts a page break and
% creates the second title. It will be ignored for other modes.
\IEEEpeerreviewmaketitle


% !TEX root =  root.tex

%\section{INTRODUCTION}
%\label{sec:introduction}

%Ezio: We need to highlight the guided synthesis more the detection
%Ezio: We need to say that can be applied to other networks of distributed systems with grid topology
%Ezio: We need to avoid widows, that means one or two words per lines

%\IEEEraisesectionheading{
\section{Introduction}\label{sec:introduction}
%}
Enormous improvements in the performance of \emph{Graphics Processing Units} (GPUs) driven by the video game market have made GPUs interesting for scientific computations. \\
The massively parallel paradigm of such processors poses new challenges on the software developers in order to make use of all the hardware resources. The Khronos Group specified the OpenCL programming language \footnote{\url{https://www.khronos.org/opencl/}} to deal with the parallelism. With the \emph{Compute Unified Device Architecture} \footnote{\url{https://developer.nvidia.com/cuda-zone}} (short CUDA) Nvidia launched its own programming language and environment for its Geforce GPUs and Tesla computation cards. \\
In this paper we will use the problem of matrix inversion as an example  for optimizing algorithms for CUDA.
\vspace{0.3cm}\\
The rest of the paper is organized as follows.  In Section~\ref{sec:problem} we will state the problem of inverting a matrix. In Section~\ref{sec:methods} we are going to describe our approach of implementing and verifying the algorithm for CUDA.\\
In Section~\ref{sec:results} we compare the performance of the parallel algorithms running on various of CUDA devices to a serial version of the algorithms running on modern x64 instruction set CPUs. Finally we will draw our conclusion in Section~\ref{sec:conclusion}.





%% !TEX root =  root.tex


\section{PROBLEM STATEMENT}\label{sec:system}




% !TEX root =  root.tex


\section{PROBLEM STATEMENT}\label{sec:problem}
First we will repeat some mathematical definitions for the sake of completeness.
\begin{definition}
	\textbf{Inverse Matrix}
	Let $A \in \mathbb{R}^n \times \mathbb{R}^n$ and $I_n$ be the identity matrix of dimension $n$, then the matrix $A^{-1} \in \mathbb{R}^n \times \mathbb{R}^n$ is called \textbf{inverse matrix} of $A$ if and only if
	\begin{equation}\label{eq:inverse}
	A \cdot A^{-1} = I_n
	\end{equation}
\end{definition}
A quadratic matrix $A$ is called \emph{regular} if there exists a inverse matrix of $A$. If there exists no inverse of $A$ then $A$ is called \emph{singular}.\\
Inverse matrices can be though to find, but have various of different uses especially in control theory to describe linear systems.  
\vspace{0.2cm}\\
There are a lot of ways to compute the inverse of a matrix, we will introduce some of the most common ones.\\
\subsection*{Gaussian-Elimination}
Calculcating the inverse by Gaussian-elemination is done by starting with two matrixes, the matrix A that should be inverted and matrix B which is initialized to the identity matrix. By performing row reductions on matrix A to make it into the identity matrix and copying all operations to matrix B gives that B is the inverse of A when A is the identity matrix.

The operations of transforming A into the identity matrix is divided into steps, forward elimination and backward elimination.
While doing forward elimination, row operations is performed on the matrix so that the diagonal is 1 and that every element below is zero.
Backward elimination is performing row operations so that all elements above the diagonal is zero.

\subsection*{Inversion by adjugate matrix}
The inverse $A^{-1}$ of $A$ can be computed by:
\begin{equation}
	A^{-1}_{i,j} := \frac{1}{\text{det}(A)} \text{adj}_{i,j}(A)
\end{equation}
with
\begin{definition}
	\textbf{Adjugate matrix} Let $M_{i,j}$ be the matrix $A$ where the $i$-th row and the $j$-th column are removed. Then:
	\[ \text{\textnormal{adj}}_{i,j}(A) = (-1)^{i+j} \text{\textnormal{det}}(M_{j,i}) \]
\end{definition}
This way of inverting a matrix works well for small $n$ (especially if $n \leq 3$), because the determinant is easy to compute in these cases.\\
\subsection*{LU decomposition}
The LU decomposition, namely with Crout's method, computes two matrices $L$ and $U$, such that $A=L \cdot U$. Where $L$ is a lower triangular matrix and $U$ is a upper triangular with 1's in the diagonal. An example with dimension 3 is:
\begin{align*}
	 \begin{pmatrix}
	 a_{11} & a_{12} & a_{13}\\
	 a_{21} & a_{22} & a_{23}\\
	 a_{31} & a_{32} & a_{33}\\
	 \end{pmatrix} & = \begin{pmatrix} l_{11}  & 0 & 0\\
	 l_{21} & l_{22} & 0\\
	 l_{31} & l_{32} & l_{33}
	 \end{pmatrix} \cdot \begin{pmatrix} 1 & u_{12} & u_{13}\\
	 0 & 1 & u_{23}\\
	 0 & 0 & 1\\
	 \end{pmatrix}\\
	& = \begin{pmatrix} l_{11} & l_{11} u_{12} & l_{11} u_{13}\\
	l_{21} & l_{21} u_{12} + l_{22} & l_{21} u_{13} + l_{22} u_{23}\\
	l_{31} & l_{31} u_{12} + l_{32} & l_{31} u_{13} + l_{32} u_{23} + l_{33}\\
	\end{pmatrix}
\end{align*} 
From this equation we will derive a recursive algorithm for computing the $L$ and $U$ matrices:
First we will set $l_{11} = a_{11}$, next we will dived the first row of $A$ by $l_{11}$, thus giving us $u_{12}$ and $u_{13}$, i.e. the first row of $U$.\\
Next we will set the remaining first column of $L$ to the first column of $A$, i.e. $l_{k,1} = a_{k,1}$. Finally we will update the right bottom 2 by 2 sub matrix by subtracting the product of the corresponding elements from the first row. I.e. the remaining 2-by-2 matrix will be of the form: 
\begin{equation*}
 \begin{pmatrix}
l_{22} & l_{22} u_{23}\\
l_{32} & l{32} u_{23} + l_{32}\\
\end{pmatrix}
\end{equation*}
Thus we can apply the same algorithm recursively on this sub-matrix.
\vspace{0.3cm}\\
During this process we divided by $l_{11}$, so in the scenario where $l_{11}=0$, it is necessary to swap the first row with a different row where the first element is not zero. Formally defined this corresponds to a row-permutation of the matrix $A$. So in fact we are looking for the two matrices $L$, $U$ and a row-permutation $\pi$, such that $\pi(A) = L \cdot U$\\  
If there is no such row, then the matrix $A$ is singular, i.e. there exists to LU decomposition and no inverse. \\
For keeping the numerical error done by this process low, it is recommended to find a row-permutation $\pi$, such that in each step the element $l_{11}$ is maximal with respect to the first column.
\vspace{0.3cm}\\
To get $A^{-1}$, every column of the inverse is computed separately, i.e. we need to find column vectors $x_i$ such that $L \cdot U \cdot x = b_i$, where $b_i$ is the i-th column of the identity matrix. Note that if we use a row-permutation $\pi$, then $\pi$ must also be applied to the vector $b_i$. The computation is again separated into two steps:\\
\begin{enumerate}
	\item Find vectors $y_i$ such that $L \cdot y_i = b_i$, this is called forward substitution
	\item Find vectors $x_i$ such that $U \cdot x_i = y_i$, this is called backward substitution
\end{enumerate}
With our example:
\begin{align*}
	\begin{pmatrix} l_{11}  & 0 & 0\\
	l_{21} & l_{22} & 0\\
	l_{31} & l_{32} & l_{33}
	\end{pmatrix} \cdot \begin{pmatrix} y_{i,1} \\ y_{i,2} \\ y_{i,3}	\end{pmatrix}
	& = \begin{pmatrix}
	b_{i,1} \\ b_{i,2} \\ b_{i,3}
	\end{pmatrix} \\
	\begin{pmatrix}
	l_{11} y_{i,1}\\ l_{21} y_{i,1} + l_{22} y_{i,2}\\
	l_{31} y_{i,1} + l_{32} y_{i,2} + l_{33} y_{i,3}
	\end{pmatrix}	& = \begin{pmatrix}
	b_{i,1} \\ b_{i,2} \\ b_{i,3}
	\end{pmatrix} \\
\end{align*}
Where we will start to compute the top most element $y_{i,1}$ first and iteratively computing the next element down to the bottom most element, this is why this process is called forward substitution.
\vspace{0.3cm}\\
Having computed $y_i$ we will next compute $x_i$, e.g.
\begin{align*}
\begin{pmatrix} 1 & u_{12} & u_{13}\\
0 & 1 & u_{23}\\
0 & 0 & 1
\end{pmatrix} \cdot \begin{pmatrix} x_{i,1} \\ x_{i,2} \\ x_{i,3}	\end{pmatrix}
& = \begin{pmatrix} y_{i,1} \\ y_{i,2} \\ y_{i,3}	\end{pmatrix}\\
\begin{pmatrix}
x_{i,1} + x_{i,2} u_{12} + x_{i,3} u_{13}\\ x_{i,2} + x_{i,3} u_{23}\\
x_{i,3}
\end{pmatrix}	& = \begin{pmatrix} y_{i,1} \\ y_{i,2} \\ y_{i,3}	\end{pmatrix}\\
\end{align*}
As the name backwards substitution reveals, this time we will start computing the bottom most element $x_{i,3}$ first and iteratively use the results to compute the second bottom most element and so on.

\subsection*{CUDA terminology}
In the following the term \emph{host} refers to program code that is executed on the CPU. On the other hand is the term \emph{device} used to refer to program code that is executed on the GPU or computation card.\\
A \emph{kernel} is a function that can be executed on the device.


% !TEX root =  root.tex


\section{METHODS}
\label{sec:methods}
\subsection*{Verification}
To verify if a matrix $A$ was inverted correctly one has to check if equation (\ref{eq:inverse}) holds. So our implementation also requires functions to multiply two matrices and check if a matrix is equal to the identity matrix. Ideally these procedures are also done using CUDA.\\
Lastly we need test matrices to run our implementation. One solution might be to read them from a file, another would be to generate the matrices randomly. We choose the latest because this way we can make use of CUDA again, in form of the cuRAND \footnote{\url{https://developer.nvidia.com/curand}} library for random number generation. So all in all our implementation will consists of the steps listed in figure (\ref{fig:steps})

\begin{figure}
\begin{tikzpicture}
	\node (entry) at (0,-0.5) [draw=none,rectangle]  {};
	\node (curand) at (0,-2) [draw,rectangle,fill=blue!10,align=center,minimum width=5cm, minimum height=1.2cm] {Generate random matrix\\using cuRAND};
	\node (inversion) at (0,-4) [draw,rectangle,fill=blue!10,align=center,minimum width=5cm, minimum height=1.2cm] {Invert matrix};
	\node (multiply) at (0,-6) [draw,rectangle,fill=blue!10,align=center,minimum width=5cm, minimum height=1.2cm] {Multiply $A$ with $A^{-1}$\\using shared memory algorithm};
	\node (check) at (0,-8) [draw,rectangle,fill=blue!10,align=center,minimum width=5cm, minimum height=1.2cm] {Check if equal to identity matrix\\using reduction} ;
	\draw[-triangle 60] (entry) to (curand);
	\draw[-triangle 60] (curand) to (inversion);
	\draw[-triangle 60] (inversion) to (multiply);
	\draw[-triangle 60] (multiply) to (check);
	
\end{tikzpicture}
\centering
\label{fig:steps}
\caption{Steps done by our implementation}
\end{figure}
\subsection*{Random matrix generation}
The cuRand library for generating pseudo and quasi random numbers consists of two major modules:
the host API and the device API.\\
The device API, which can be accessed by \texttt{curand\_kernel.h}, are low level functions that can only be called from within a kernel. \\
The host API on the other hand, can be access by \texttt{curand.h} and contains high level functions for host code. With the host API, entire arrays can be filled with random numbers using only one API call. \\
Because a row-major-ordering matrix of dimension $n$ by $n$ is represented as an array of size $n^2$, a random matrix can be generated with only one API call. 

\subsection*{Matrix multiplication using shared memory}
For the multiplication of $AA^{-1}$ we used the shared memory algorithm for matrix multiplication found in the CUDA 7.5 samples directory, namely the kernel in \texttt{0\_Simple/matrixMul\_nvrtc/matrixMul\_kernel.cu}.\\
The algorithms works like visualized in figure (\ref{fig:multiply}):
The blue part of $A \cdot B$ is computed by exactly one warp. Each entry of $A \cdot B$ is computed by one thread.\\ 
In each step the green sub parts of $A$ and $B$ are loaded into the shared memory cache, so the threads of the warp can access them quickly. The threads multiply the columns of $A$ that are in the shared memory with the rows of $B$ that are in the shared memory and add the result to the corresponding entries in $A \cdot B$.
\vspace{0.3cm}\\
The original algorithms from the samples directory does not take matrices into account where the dimensions are not a multiple to the predefined shared memory block size. So we added the support for the multiplication of quadratic matrices of arbitrary dimension. This was done, by initializing the shared memory cache with 0, for those elements loaded into the cache that would exceed the dimension of the matrix.
\begin{figure}
\begin{tikzpicture}
	\node at (-0.85,1.8) {$A$};
	\node at (1.8,2.25) {$B$};
	\node at (2.0,0.75) {$A \cdot B$};
	\node at (0.75,0.75) [rectangle,fill=blue!40,minimum width=0.5cm, minimum height=0.5cm] {};
	\node at (-1.35,0.75) [rectangle,fill=green!60,minimum width=0.5cm, minimum height=0.5cm] {};
	\node at (0.75,2.85) [rectangle,fill=green!60,minimum width=0.5cm, minimum height=0.5cm] {};
	\draw[step=0.5,black] grid (1.5,1.5);
	\draw[step=0.5,black,xshift=-1.6cm] grid (1.5,1.5);
	\draw[step=0.5,black,yshift=1.6cm] grid (1.5,1.5);
	\node at (-2.6,2.25) {\textbf{Step 1}};
	
	\draw[dashed] (-3.0,-0.25) to (3,-0.25);
	\draw[dashed] (-3.0,-4) to (3,-4);
		\node at (-0.85,-2) {$A$};
		\node at (1.8,-1.45) {$B$};
		\node at (2.0,-2.95) {$A \cdot B$};
		\node at (0.75,-2.95) [rectangle,fill=blue!40,minimum width=0.5cm, minimum height=0.5cm] {};
		\node at (-0.85,-2.95) [rectangle,fill=green!60,minimum width=0.5cm, minimum height=0.5cm] {};
		\node at (0.75,-1.35) [rectangle,fill=green!60,minimum width=0.5cm, minimum height=0.5cm] {};
		\draw[step=0.5,black,yshift=-3.7cm] grid (1.5,1.5);
		\draw[step=0.5,black,xshift=-1.6cm,yshift=-3.7cm] grid (1.5,1.5);
		\draw[step=0.5,black,yshift=-2.1cm] grid (1.5,1.5);
		\node at (-2.6,-1.3) {\textbf{Step 2}};
		
		\node at (-0.85,-5.8) {$A$};
		\node at (1.8,-5.25) {$B$};
		\node at (2.0,-6.75) {$A \cdot B$};
		\node at (0.75,-6.75) [rectangle,fill=blue!40,minimum width=0.5cm, minimum height=0.5cm] {};
		\node at (-0.35,-6.75) [rectangle,fill=green!60,minimum width=0.5cm, minimum height=0.5cm] {};
		\node at (0.75,-5.65) [rectangle,fill=green!60,minimum width=0.5cm, minimum height=0.5cm] {};
		\draw[step=0.5,black,yshift=-7.5cm] grid (1.5,1.5);
		\draw[step=0.5,black,xshift=-1.6cm,yshift=-7.5cm] grid (1.5,1.5);
		\draw[step=0.5,black,yshift=-5.9cm] grid (1.5,1.5);
		\node at (-2.6,-5.2) {\textbf{Step 3}};
		
\end{tikzpicture}
\centering
\label{fig:multiply}
\caption{Multiplication of $A$ and $B$ using shared memory}
\end{figure}

\subsection*{Checking if a matrix is the identity matrix}
We need a procedure which checks if a given matrix $A$ of dimension $n$-by-$n$ is equal to the identity matrix $I_n$. Formally defined:
\begin{equation*}
is\_id(A) := \begin{cases}
1 & \text{if } A \text{ has 1's in the diagonal, and 0's elsewhere}\\
0 & \text{otherwise}
\end{cases}
\end{equation*}
So $is\_id$ maps a matrix of $\mathbb{R}^{n \times n}$ to a single value $\{0,1 \}$. Therefore a implementation of $is\_id$ can be done efficiently using a \textbf{reduction}.  \\
We adapted the code from the CUDA 7.5 samples directory, namely from \texttt{/samples/6\_Advanced/reduction/}. The sample code is based on the whitepaper from Harris \cite{Harris2010} and optimizes the reduction for a summation of an array. The modification we did was that instead of summing up all array entries we used a auxiliary function $f_{acc}$ which takes an entry and its index as input and is returns 1 if  the value and the index to not match with the definition of the identity matrix. Formally:
\begin{equation}
f_{acc}(A_{i,j}) := \begin{cases} 1 & \text{if }i = j \text{ and } |A_{i,j}-1| \geq \varepsilon\\
1 & \text{if } i \neq j \text{ and } |A_{i,j}| \geq \varepsilon\\
0 & \text{otherwise}	\end{cases}
\end{equation}
Note that due to the limited floating point representation, a rounding error variable $\varepsilon$ must be chosen that should be close to 0, e.g. $\varepsilon = 10^{-6}$. \\
Our modified reduction computes $\sum_{i=0,j=0}^{n,n} f_{acc}(A_{i,j})$, so if the matrix $A$ matches the definition of the identity matrix the sum will be 0, if at least one entry of the matrix does not match the definition of the identity matrix, then the sum will be greater than 0.

\subsection*{LU decomposition}
To store the L and the U matrix, only one array of size $n^2$ is needed, because of the fixed 0s and 1s. For example the matrices
\begin{equation*}
 L = \begin{pmatrix} l_{11}  & 0 & 0\\ 
l_{21} & l_{22} & 0\\
l_{31} & l_{32} & l_{33}
\end{pmatrix}, U = \begin{pmatrix} 1 & u_{12} & u_{13}\\
0 & 1 & u_{23}\\
0 & 0 & 1\\
\end{pmatrix} 
\end{equation*}
Can be stored as
\begin{equation*}
LU = \begin{pmatrix} l_{11}  &  u_{12} & u_{13}\\
l_{21} & l_{22} & u_{23}\\
l_{31} & l_{32} & l_{33}
\end{pmatrix}
\end{equation*}
and thereby reducing the memory requirements by half.
\vspace{0.3cm}\\
The computation of the LU matrices and the row-permutation is separated into the following steps:
\begin{enumerate}
	\item Find Pivot, i.e. find row with the maximal absolute value in first column.\\
	This can be done efficiently on the GPU using a \textbf{reduction}. Like with the check for identity matrix we used the sample reduction in \texttt{/samples/6\_Advanced/reduction/} by Harris \cite{Harris2010}.
	\item Divide the pivot-row by pivot. Because the single computations do not depend on each other one grid equal to the side of the row can be launched. Each thread of the grid divides exactly one element of the row by the pivot. 
	\item Swap rows. This is application the row-permutation. Like before, each thread will swap exactly one item of the pivot row with the item in the first row.
	\item Update sub-matrix. As mentioned in section (\ref{sec:problem}) each element in the (n-1)-by-(n-1) sub-matrix needs to be subtracted by the product of the corresponding elements in the first row and first column. Every entry of this sub-matrix is updated by one thread. So the grid launched to do this task if of size (n-1)-by-(n-1)
	\item Process sub-matrix (recursive step)
\end{enumerate}
The last step would involve a recursive call to the procedure to itself. For performance reasons, the recursive algorithm was transformed into an iterative algorithm.
\subsection*{Matrix inversion of LU matrix}
As stated in section (\ref{sec:problem}), each column of the inverse matrix is computed separately by solving a linear equation. Because each linear equation is independent from the others, they can be solved in parallel. So a grid of size $n$ is launched, where each thread solves exactly one linear equation.  

% !TEX root =  root.tex

\section{RESULTS}\label{sec:results}

%% !TEX root = root.tex

\section{RELATED WORK}
\label{sec:related}



% !TEX root =  root.tex

\section{CONCLUSION}
\label{sec:conclusion}

GPGPU computing can be used to speed up matrix inversion for medium to large matrices, for small matrices the overhead of passing data back and forth between the host and the GPU is bigger than just running the inversion on the CPU.
\FloatBarrier

% Computer Society journal (but not conference!) papers do something unusual
% with the very first section heading (almost always called "Introduction").
% They place it ABOVE the main text! IEEEtran.cls does not automatically do
% this for you, but you can achieve this effect with the provided
% \IEEEraisesectionheading{} command. Note the need to keep any \label that
% is to refer to the section immediately after \section in the above as
% \IEEEraisesectionheading puts \section within a raised box.




% The very first letter is a 2 line initial drop letter followed
% by the rest of the first word in caps (small caps for compsoc).
% 
% form to use if the first word consists of a single letter:
% \IEEEPARstart{A}{demo} file is ....
% 
% form to use if you need the single drop letter followed by
% normal text (unknown if ever used by IEEE):
% \IEEEPARstart{A}{}demo file is ....
% 
% Some journals put the first two words in caps:
% \IEEEPARstart{T}{his demo} file is ....
% 
% Here we have the typical use of a "T" for an initial drop letter
% and "HIS" in caps to complete the first word.
%\IEEEPARstart{T}{his} demo file is intended to serve as a ``starter file''
%for IEEE Computer Society journal papers produced under \LaTeX\ using
%IEEEtran.cls version 1.8a and later.
% You must have at least 2 lines in the paragraph with the drop letter
% (should never be an issue)
%I wish you the best of success.

%\hfill mds
 
%\hfill September 17, 2014

%\subsection{Subsection Heading Here}
%Subsection text here.

% needed in second column of first page if using \IEEEpubid
%\IEEEpubidadjcol

%\subsubsection{Subsubsection Heading Here}
%Subsubsection text here.


% An example of a floating figure using the graphicx package.
% Note that \label must occur AFTER (or within) \caption.
% For figures, \caption should occur after the \includegraphics.
% Note that IEEEtran v1.7 and later has special internal code that
% is designed to preserve the operation of \label within \caption
% even when the captionsoff option is in effect. However, because
% of issues like this, it may be the safest practice to put all your
% \label just after \caption rather than within \caption{}.
%
% Reminder: the "draftcls" or "draftclsnofoot", not "draft", class
% option should be used if it is desired that the figures are to be
% displayed while in draft mode.
%
%\begin{figure}[!t]
%\centering
%\includegraphics[width=2.5in]{myfigure}
% where an .eps filename suffix will be assumed under latex, 
% and a .pdf suffix will be assumed for pdflatex; or what has been declared
% via \DeclareGraphicsExtensions.
%\caption{Simulation results for the network.}
%\label{fig_sim}
%\end{figure}

% Note that IEEE typically puts floats only at the top, even when this
% results in a large percentage of a column being occupied by floats.
% However, the Computer Society has been known to put floats at the bottom.


% An example of a double column floating figure using two subfigures.
% (The subfig.sty package must be loaded for this to work.)
% The subfigure \label commands are set within each subfloat command,
% and the \label for the overall figure must come after \caption.
% \hfil is used as a separator to get equal spacing.
% Watch out that the combined width of all the subfigures on a 
% line do not exceed the text width or a line break will occur.
%
%\begin{figure*}[!t]
%\centering
%\subfloat[Case I]{\includegraphics[width=2.5in]{box}%
%\label{fig_first_case}}
%\hfil
%\subfloat[Case II]{\includegraphics[width=2.5in]{box}%
%\label{fig_second_case}}
%\caption{Simulation results for the network.}
%\label{fig_sim}
%\end{figure*}
%
% Note that often IEEE papers with subfigures do not employ subfigure
% captions (using the optional argument to \subfloat[]), but instead will
% reference/describe all of them (a), (b), etc., within the main caption.
% Be aware that for subfig.sty to generate the (a), (b), etc., subfigure
% labels, the optional argument to \subfloat must be present. If a
% subcaption is not desired, just leave its contents blank,
% e.g., \subfloat[].


% An example of a floating table. Note that, for IEEE style tables, the
% \caption command should come BEFORE the table and, given that table
% captions serve much like titles, are usually capitalized except for words
% such as a, an, and, as, at, but, by, for, in, nor, of, on, or, the, to
% and up, which are usually not capitalized unless they are the first or
% last word of the caption. Table text will default to \footnotesize as
% IEEE normally uses this smaller font for tables.
% The \label must come after \caption as always.
%
%\begin{table}[!t]
%% increase table row spacing, adjust to taste
%\renewcommand{\arraystretch}{1.3}
% if using array.sty, it might be a good idea to tweak the value of
% \extrarowheight as needed to properly center the text within the cells
%\caption{An Example of a Table}
%\label{table_example}
%\centering
%% Some packages, such as MDW tools, offer better commands for making tables
%% than the plain LaTeX2e tabular which is used here.
%\begin{tabular}{|c||c|}
%\hline
%One & Two\\
%\hline
%Three & Four\\
%\hline
%\end{tabular}
%\end{table}


% Note that the IEEE does not put floats in the very first column
% - or typically anywhere on the first page for that matter. Also,
% in-text middle ("here") positioning is typically not used, but it
% is allowed and encouraged for Computer Society conferences (but
% not Computer Society journals). Most IEEE journals/conferences use
% top floats exclusively. 
% Note that, LaTeX2e, unlike IEEE journals/conferences, places
% footnotes above bottom floats. This can be corrected via the
% \fnbelowfloat command of the stfloats package.




%\section{Conclusion}
%The conclusion goes here.





% if have a single appendix:
%\appendix[Proof of the Zonklar Equations]
% or
%\appendix  % for no appendix heading
% do not use \section anymore after \appendix, only \section*
% is possibly needed

% use appendices with more than one appendix
% then use \section to start each appendix
% you must declare a \section before using any
% \subsection or using \label (\appendices by itself
% starts a section numbered zero.)
%


%\appendices
%\section{Proof of the First Zonklar Equation}
%Appendix one text goes here.

% you can choose not to have a title for an appendix
% if you want by leaving the argument blank
%\section{}
%Appendix two text goes here.


% Can use something like this to put references on a page
% by themselves when using endfloat and the captionsoff option.
\ifCLASSOPTIONcaptionsoff
  \newpage
\fi



% trigger a \newpage just before the given reference
% number - used to balance the columns on the last page
% adjust value as needed - may need to be readjusted if
% the document is modified later
%\IEEEtriggeratref{8}
% The "triggered" command can be changed if desired:
%\IEEEtriggercmd{\enlargethispage{-5in}}

% references section

% can use a bibliography generated by BibTeX as a .bbl file
% BibTeX documentation can be easily obtained at:
% http://www.ctan.org/tex-archive/biblio/bibtex/contrib/doc/
% The IEEEtran BibTeX style support page is at:
% http://www.michaelshell.org/tex/ieeetran/bibtex/
%\bibliographystyle{IEEEtran}
% argument is your BibTeX string definitions and bibliography database(s)
%\bibliography{IEEEabrv,../bib/paper}
%
% <OR> manually copy in the resultant .bbl file
% set second argument of \begin to the number of references
% (used to reserve space for the reference number labels box)
%\begin{thebibliography}{1}

%\bibitem{IEEEhowto:kopka}
%H.~Kopka and P.~W. Daly, \emph{A Guide to \LaTeX}, 3rd~ed.\hskip 1em plus
 % 0.5em minus 0.4em\relax Harlow, England: Addison-Wesley, 1999.

%\end{thebibliography}

% biography section
% 
% If you have an EPS/PDF photo (graphicx package needed) extra braces are
% needed around the contents of the optional argument to biography to prevent
% the LaTeX parser from getting confused when it sees the complicated
% \includegraphics command within an optional argument. (You could create
% your own custom macro containing the \includegraphics command to make things
% simpler here.)
%\begin{IEEEbiography}[{\includegraphics[width=1in,height=1.25in,clip,keepaspectratio]{mshell}}]{Michael Shell}
% or if you just want to reserve a space for a photo:

%\begin{IEEEbiography}{Michael Shell}
%Biography text here.
%\end{IEEEbiography}

% if you will not have a photo at all:
%\begin{IEEEbiographynophoto}{John Doe}
%Biography text here.
%\end{IEEEbiographynophoto}

% insert where needed to balance the two columns on the last page with
% biographies
%\newpage

%\begin{IEEEbiographynophoto}{Jane Doe}
%Biography text here.
%\end{IEEEbiographynophoto}

% You can push biographies down or up by placing
% a \vfill before or after them. The appropriate
% use of \vfill depends on what kind of text is
% on the last page and whether or not the columns
% are being equalized.

%\vfill

% Can be used to pull up biographies so that the bottom of the last one
% is flush with the other column.
%\enlargethispage{-5in}

%\printbibliography 

% that's all folks

\bibliography{nfm}
%\begin{thebibliography}{99}
%
%\bibitem{c1}
%J.G.F. Francis, The QR Transformation I, {\it Comput. J.}, vol. 4, 1961, pp 265-271.
%
%\bibitem{c2}
%H. Kwakernaak and R. Sivan, {\it Modern Signals and Systems}, Prentice Hall, Englewood Cliffs, NJ; 1991.
%
%\bibitem{c3}
%D. Boley and R. Maier, "A Parallel QR Algorithm for the Non-Symmetric Eigenvalue Algorithm", {\it in Third SIAM Conference on Applied Linear Algebra}, Madison, WI, 1988, pp. A20.
%
%\end{thebibliography}
\end{document}


